% ARCHIVO GENERADO por generate_software_doc.py -- no editar a mano.
\chapter{Software}
\label{chap:software}
\lstset{language=Python,basicstyle=\ttfamily\scriptsize,numbers=left,numberstyle=\tiny,breaklines=true,frame=single,keywordstyle=\color{blue},commentstyle=\color{gray},inputencoding=utf8,extendedchars=true,literate={á}{{\'a}}1 {é}{{\'e}}1 {í}{{\'i}}1 {ó}{{\'o}}1 {ú}{{\'u}}1 {ñ}{{\~n}}1 {Á}{{\'A}}1 {É}{{\'E}}1 {Í}{{\'I}}1 {Ó}{{\'O}}1 {Ú}{{\'U}}1 {Ñ}{{\~N}}1}
Este capítulo documenta el software del robot Storm Breaker (MakeX 2026). Se genera automáticamente a partir de los comentarios Doxygen de Robot.py.

\section{Descripci\'on general}
Control principal del robot con movimiento omnidireccional, griper, lanzador y visión QR.

Este módulo integra todos los subsistemas del robot: - Tracción omnidireccional de 4 ruedas mediante motores encoder (M1–M4). - Griper con movimiento vertical (DC1) y apertura/cierre de pinzas (SM1). - Lanzador de pelotas con dos motores brushless (BL1, BL2). - Visión artificial con cámara de IA para lectura y comparación de etiquetas QR. - Control manual mediante gamepad con mapeo de teclas a acciones. - Modo autónomo de 28 s que usa la cámara para seleccionar y alcanzar el cubo de mayor valor detectado en campo.

Archivo fuente: \texttt{Robot.py} (540 l\'ineas).
\section{Hardware y variables globales}
\begin{longtable}{@{}p{0.28\linewidth}p{0.66\linewidth}@{}}\toprule
\textbf{Nombre} & \textbf{Descripci\'on} \\\midrule\endhead
\texttt{pot\_lanzamiento} & Potencia del lanzador (porcentaje) aplicada a BL1/BL2. \\
\texttt{aper\_pot\_griper} & Potencia del motor de apertura/cierre de la garra (SM1). \\
\texttt{mov\_pot\_griper} & Potencia del motor de movimiento de la garra (DC1). \\
\texttt{v1} & Velocidad objetivo de la rueda M1 (la calcula Checar\_Velocidades). \\
\texttt{v2} & Velocidad objetivo de la rueda M2. \\
\texttt{v3} & Velocidad objetivo de la rueda M3. \\
\texttt{v4} & Velocidad objetivo de la rueda M4. \\
\texttt{Bandera} & Bandera para salir de ciclos (1 = terminado). \\
\texttt{TECLAS} & Tupla de nombres de teclas del gamepad que están registradas en el sistema. \\
\texttt{acciones\_gamepad} & Diccionario que mapea cada tecla registrada a su función y argumentos. \\
\texttt{TECLAS} &  \\
\texttt{acciones\_gamepad} &  \\
\bottomrule\end{longtable}
\section{C\'omo se ejecuta el programa}
\subsection{Secuencia de arranque}
Instrucciones a nivel de m\'odulo, ejecutadas en orden al encender el NovaPi:

\begin{figure}[H]\centering
\begin{tikzpicture}[>=Stealth,
  st/.style={draw,rounded corners,fill=green!10,align=center,minimum width=3.2cm,minimum height=0.8cm,font=\footnotesize},
  fin/.style={draw,ellipse,fill=gray!15,font=\footnotesize}]
\node[fin] (start) at (0,0) {Encendido};
\node[st,below=0.6cm of start] (b0) {\texttt{sleep(30)}};
\draw[->] (start) -- (b0);
\node[st,below=0.6cm of b0] (b1) {\texttt{InicializarTag()}};
\draw[->] (b0) -- (b1);
\node[st,below=0.6cm of b1] (b2) {\texttt{main()}};
\draw[->] (b1) -- (b2);
\end{tikzpicture}
\caption{Secuencia de arranque (generada del c\'odigo).}\end{figure}
\subsection{Ciclo principal y botones del gamepad}
\subsection{Flujo de decisi\'on aut\'onomo}
% Diagrama hecho a mano de Comparacion_QR(); actualizar si cambia la lógica.
\begin{figure}[H]\centering
\resizebox{\linewidth}{!}{%
\begin{tikzpicture}[>=Stealth,
  st/.style={draw,rounded corners,fill=orange!10,align=center,minimum width=3.4cm,font=\footnotesize},
  dc/.style={draw,diamond,aspect=2.4,fill=yellow!15,align=center,font=\scriptsize,inner sep=1pt}]
\node[st] (read) at (0,0) {Leer tag izquierdo (id 2)\\y derecho (id 3)};
\node[dc,below=0.6cm of read] (none) {¿Ninguno\\detectado?};
\node[st,right=1.5cm of none] (cones) {Dos conos:\\\texttt{Mov\_Auto(1, 55)}};
\node[dc,below=0.6cm of none] (same) {¿Misma posici\'on\\en ambos tags?};
\node[dc,below=0.6cm of same] (cube) {¿Tag central\\= 2?};
\node[st,right=1.5cm of cube] (grab) {Agarrar cubo};
\node[st,below=0.6cm of cube] (move) {Avanzar \texttt{Mov\_Auto(1, 44)},\\releer tag y comparar};
\node[st,right=1.5cm of same] (cmp) {Comparar ambos tags\\\texttt{Comparacion\_Valores(..., 0)}};
\node[st,below=0.6cm of move] (val) {\texttt{Comparacion\_Valores(..., 1)}\\\texttt{Mov\_Auto}};
\draw[->] (read) -- (none);
\draw[->] (none) -- node[above,font=\scriptsize]{s\'i} (cones);
\draw[->] (none) -- node[right,font=\scriptsize]{no} (same);
\draw[->] (same) -- node[right,font=\scriptsize]{s\'i} (cube);
\draw[->] (same) -- node[above,font=\scriptsize]{no} (cmp);
\draw[->] (cube) -- node[above,font=\scriptsize]{s\'i} (grab);
\draw[->] (cube) -- node[right,font=\scriptsize]{no} (move);
\draw[->] (move) -- (val);
\end{tikzpicture}}
\caption{Flujo de decisi\'on de \texttt{Comparacion\_QR()} (se repite durante 28 s).}
\end{figure}

\section{Arquitectura: llamadas entre funciones}
\begin{figure}[H]\centering\resizebox{\linewidth}{!}{%
\begin{tikzpicture}[node distance=1cm, >=Stealth,
  fn/.style={draw,rounded corners,fill=blue!8,font=\footnotesize\ttfamily,minimum height=0.7cm,inner sep=4pt}]
\node[fn] (InicializarTag) at (0.00,0.00) {InicializarTag};
\node[fn] (MovGriper) at (0.00,-1.10) {Mov\_Griper};
\node[fn] (AperturaGriper) at (0.00,-2.20) {Apertura\_Griper};
\node[fn] (AnguloBrush) at (0.00,-3.30) {Angulo\_Brush};
\node[fn] (modoautomatico) at (0.00,-4.40) {modo\_automatico};
\node[fn] (lanzamiento) at (0.00,-5.50) {lanzamiento};
\node[fn] (ResetMotores) at (0.00,-6.60) {Reset\_Motores};
\node[fn] (main) at (0.00,-7.70) {main};
\node[fn] (ComparacionQR) at (5.20,0.00) {Comparacion\_QR};
\node[fn] (ChecarServo) at (5.20,-1.10) {ChecarServo};
\node[fn] (Lanzador) at (5.20,-2.20) {Lanzador};
\node[fn] (getpressedkey) at (5.20,-3.30) {get\_pressed\_key};
\node[fn] (ComparacionValores) at (10.40,0.00) {Comparacion\_Valores};
\node[fn] (MovAuto) at (15.60,0.00) {Mov\_Auto};
\node[fn] (ChecarVelocidades) at (20.80,0.00) {Checar\_Velocidades};
\node[fn] (VelMotores) at (20.80,-1.10) {Vel\_Motores};
\draw[->] (ComparacionQR) -- (MovAuto);
\draw[->] (ComparacionQR) -- (ComparacionValores);
\draw[->] (MovAuto) -- (ChecarVelocidades);
\draw[->] (MovAuto) -- (VelMotores);
\draw[->] (ComparacionValores) -- (MovAuto);
\draw[->] (AperturaGriper) -- (ChecarServo);
\draw[->] (modoautomatico) -- (ComparacionQR);
\draw[->] (lanzamiento) -- (Lanzador);
\draw[->] (main) -- (getpressedkey);
\draw[->] (main) -- (ChecarVelocidades);
\draw[->] (main) -- (VelMotores);
\end{tikzpicture}
}\caption{Grafo de llamadas (generado con el an\'alisis del c\'odigo).}\end{figure}
\section{Decisiones de dise\~no}
\medskip\noindent\textbf{SW-DEC-001 -- Un solo archivo para NovaPi}\par
\textbf{Decisi\'on:} Todo el programa vive en un único archivo Robot.py.\\
\textbf{Justificaci\'on:} El NovaPi de mBuild carga un solo script de Python; mantener todo junto evita problemas de importación y facilita el cambio rápido en competencia.

\medskip\noindent\textbf{SW-DEC-002 -- Cinemática mecanum con signos por montaje}\par
\textbf{Decisi\'on:} Checar\_Velocidades calcula v1..v4 y Vel\_Motores invierte M2 y M4.\\
\textbf{Justificaci\'on:} Separar el cálculo de la aplicación permite probar la cinemática sin mover motores; los motores del lado derecho están montados en sentido opuesto.

\medskip\noindent\textbf{SW-DEC-003 -- Control manual siempre activo}\par
\textbf{Decisi\'on:} Los joysticks mueven el chasis en cada iteración de main(), incluso después de presionar un botón.\\
\textbf{Justificaci\'on:} El conductor conserva el control del chasis y puede interrumpir las acciones de los mecanismos.

\medskip\noindent\textbf{SW-DEC-004 -- Autónomo de 28 s basado en tags}\par
\textbf{Decisi\'on:} El botón N1 ejecuta Comparacion\_QR en un ciclo limitado por novapi.timer() < 28.\\
\textbf{Justificaci\'on:} El periodo autónomo es de 30 s; se deja 2 s de margen para no invadir el periodo manual.

\medskip\noindent\textbf{SW-DEC-005 -- Conversión de píxeles a cm}\par
\textbf{Decisi\'on:} Las distancias de la cámara se dividen entre 16.84 px/cm.\\
\textbf{Justificaci\'on:} Constante obtenida de la calibración del tag de 10 cm; debe recalibrarse si cambia la cámara o su posición.

\medskip\noindent\textbf{SW-DEC-006 -- Movimiento autónomo por tiempo}\par
\textbf{Decisi\'on:} Mov\_Auto mueve el robot por un tiempo fijo en lugar de usar la distancia recibida.\\
\textbf{Justificaci\'on:} Solución provisional; la conversión de distancia a velocidad/tiempo queda pendiente de calibración.

\section{Referencia de funciones y c\'odigo}
\subsection{\texttt{Checar\_Velocidades}}
\label{fn:ChecarVelocidades}
Calcula las velocidades de los cuatro motores para un robot omnidireccional.

Aplica la ecuación de cinemática inversa holonómica para una configuración de cuatro ruedas mecanum o similares. La velocidad angular se pondera con el factor @p ri antes de sumarse o restarse a cada rueda:

\begin{tabular}{@{}p{0.15\linewidth}p{0.78\linewidth}@{}}
\textbf{Par\'ametro} & \textbf{Descripci\'on} \\
\texttt{Vx} & Componente de velocidad en el eje X del robot (avance/retroceso). Valores positivos avanzan, negativos retroceden. \\
\texttt{Vy} & Componente de velocidad en el eje Y del robot (desplazamiento lateral). Valores positivos desplazan a la derecha, negativos a la izquierda. \\
\texttt{W} & Velocidad angular de giro sobre el eje Z del robot. Valores positivos giran en sentido horario, negativos en antihorario. \\
\texttt{ri} & Factor de escala (radio efectivo) que pondera la contribución angular. Normalmente se pasa 1 para mantener las unidades consistentes. \\
\end{tabular}\\[4pt]
\lstinputlisting[firstline=86,lastline=93,firstnumber=86]{../../../software/robot/main/Robot.py}
\subsection{\texttt{Vel\_Motores}}
\label{fn:VelMotores}
Aplica las velocidades especificadas a los cuatro motores encoder del robot.

Los motores M2 y M4 reciben la velocidad multiplicada por -1 para compensar su orientación física invertida respecto a M1 y M3. De este modo, una misma velocidad lógica produce el movimiento correcto en todos los motores. Esta función no realiza cálculos; solo envía los valores a los drivers.

\begin{tabular}{@{}p{0.15\linewidth}p{0.78\linewidth}@{}}
\textbf{Par\'ametro} & \textbf{Descripci\'on} \\
\texttt{v1} & Velocidad deseada para el motor M1 (float, rango típico -100 a 100). \\
\texttt{v2} & Velocidad deseada para el motor M2 (float, rango típico -100 a 100). Se aplicará con signo invertido al motor físico. \\
\texttt{v3} & Velocidad deseada para el motor M3 (float, rango típico -100 a 100). \\
\texttt{v4} & Velocidad deseada para el motor M4 (float, rango típico -100 a 100). Se aplicará con signo invertido al motor físico. \\
\end{tabular}\\[4pt]
\lstinputlisting[firstline=108,lastline=112,firstnumber=108]{../../../software/robot/main/Robot.py}
\subsection{\texttt{Comparacion\_QR}}
\label{fn:ComparacionQR}
Lee las etiquetas QR de la cámara y decide el movimiento autónomo del robot.

Consulta los slots 2 y 3 de la cámara (cubos izquierdo y derecho) y aplica la siguiente lógica de decisión: - Si ambos slots devuelven 'Null' -> se asume presencia de 2 conos; el robot avanza 55 unidades hacia la izquierda (tipo 1). - Si las distancias de los slots 2 y 3 son iguales -> los cubos están alineados en un solo bloque; se analiza el slot 1 individualmente: - Si su valor es 2 (máximo posible), se imprime "Agarrar Cubo" y no se realiza movimiento adicional. - En caso contrario, se guarda el valor y la posición del primer cubo, el robot avanza 44 u para posicionarse frente al segundo cubo, se re-lee el QR y se llama a Comparacion\_Valores() en modo 1. - Si las distancias son distintas -> los cubos están en posiciones diferentes; se llama directamente a Comparacion\_Valores() en modo 0 con los datos de los slots 2 y 3.

Llama a: \texttt{Mov\_Auto} (\ref{fn:MovAuto}), \texttt{Comparacion\_Valores} (\ref{fn:ComparacionValores})\\
\lstinputlisting[firstline=133,lastline=171,firstnumber=133]{../../../software/robot/main/Robot.py}
\subsection{\texttt{Mov\_Auto}}
\label{fn:MovAuto}
Desplaza el robot de forma autónoma en la dirección indicada durante 20 segundos.

Calcula las velocidades mediante Checar\_Velocidades(), reinicia el temporizador interno de novapi y mantiene los motores activos hasta que el temporizador supera los 20 segundos. El parámetro @p distancia está reservado para una futura implementación que escale el tiempo de movimiento; actualmente no afecta el comportamiento. - Tipo 1: movimiento lateral izquierda  (Vy = +1). - Tipo 2: movimiento lateral derecha    (Vy = -1).

\begin{tabular}{@{}p{0.15\linewidth}p{0.78\linewidth}@{}}
\textbf{Par\'ametro} & \textbf{Descripci\'on} \\
\texttt{tipo} & Entero que indica la dirección de movimiento. 1 = desplazamiento hacia la izquierda. 2 = desplazamiento hacia la derecha. Cualquier otro valor no produce movimiento. \\
\texttt{distancia} & Distancia objetivo en centímetros (parámetro reservado, no utilizado en el cálculo actual del tiempo de movimiento). \\
\end{tabular}\\[4pt]
Llama a: \texttt{Checar\_Velocidades} (\ref{fn:ChecarVelocidades}), \texttt{Vel\_Motores} (\ref{fn:VelMotores})\\
\lstinputlisting[firstline=189,lastline=201,firstnumber=189]{../../../software/robot/main/Robot.py}
\subsection{\texttt{InicializarTag}}
\label{fn:InicializarTag}
Configura la cámara de IA para detectar y analizar etiquetas QR.

Realiza la secuencia de inicialización en tres pasos: 1. Activa el modo de reconocimiento de etiquetas (switch de función 4). 2. Espera 10 segundos para que el hardware de la cámara se estabilice. 3. Establece el modo de función de etiqueta en 3 (detección de QR/ArUco) y define el tamaño físico real de las etiquetas en 10 cm, necesario para que la cámara calcule distancias en unidades métricas.

\lstinputlisting[firstline=212,lastline=216,firstnumber=212]{../../../software/robot/main/Robot.py}
\subsection{\texttt{Comparacion\_Valores}}
\label{fn:ComparacionValores}
Compara los valores QR de dos cubos y mueve el robot hacia el de mayor valor.

Convierte primero las distancias de píxeles a centímetros dividiendo entre el factor de calibración 16.84 px/cm. Luego, dependiendo del @p modo, aplica una de dos estrategias: Modo 1 — Cubos en bloques distintos: - Si Val1 > Val2 -> el cubo 1 es el objetivo. Si está a la izquierda (Pos1 == 1) se invierte Dist1 para moverse en esa dirección. El robot se desplaza (-44 + Dist1) unidades desde el punto de referencia. - Si Val2 >= Val1 -> el cubo 2 es el objetivo. Si está a la izquierda (Pos2 == 1) se invierte Dist2. El robot se desplaza Dist2 unidades. Modo 0 — Cubos en el mismo bloque: - Si Val1 > Val2 -> el cubo 1 (izquierda) es el objetivo; Dist1 se invierte y el robot se desplaza esa distancia hacia la izquierda. - Si Val2 >= Val1 -> el cubo 2 (derecha) es el objetivo; el robot se desplaza Dist2 unidades hacia la derecha.

\begin{tabular}{@{}p{0.15\linewidth}p{0.78\linewidth}@{}}
\textbf{Par\'ametro} & \textbf{Descripci\'on} \\
\texttt{Val1} & Valor numérico leído del QR del cubo 1 (entero, mayor = mejor). \\
\texttt{Dist1} & Distancia en píxeles desde el punto de referencia de la cámara hasta el QR del cubo 1. Se convierte internamente a centímetros. \\
\texttt{Pos1} & Posición lateral del cubo 1 respecto al centro de la imagen. 1 = cubo a la izquierda (coordenada X < 320 px). 0 = cubo a la derecha  (coordenada X >= 320 px). \\
\texttt{Val2} & Valor numérico leído del QR del cubo 2 (entero, mayor = mejor). \\
\texttt{Dist2} & Distancia en píxeles desde el punto de referencia de la cámara hasta el QR del cubo 2. Se convierte internamente a centímetros. \\
\texttt{Pos2} & Posición lateral del cubo 2 respecto al centro de la imagen. 1 = cubo a la izquierda (coordenada X < 320 px). 0 = cubo a la derecha  (coordenada X >= 320 px). \\
\texttt{modo} & Estrategia de comparación a aplicar. 1 = los cubos están en bloques distintos (se usa offset de 44 u). 0 = los cubos están en el mismo bloque (desplazamiento directo). \\
\end{tabular}\\[4pt]
Llama a: \texttt{Mov\_Auto} (\ref{fn:MovAuto})\\
\lstinputlisting[firstline=253,lastline=275,firstnumber=253]{../../../software/robot/main/Robot.py}
\subsection{\texttt{Mov\_Griper}}
\label{fn:MovGriper}
Controla el movimiento vertical del griper (ascenso o descenso).

Establece la potencia de mov\_pot\_griper en 50 \% y la aplica al canal DC1 de la placa de expansión de potencia. El signo positivo hace ascender el griper y el negativo lo hace descender. La variable global mov\_pot\_griper se actualiza en cada llamada.

\begin{tabular}{@{}p{0.15\linewidth}p{0.78\linewidth}@{}}
\textbf{Par\'ametro} & \textbf{Descripci\'on} \\
\texttt{modo} & Entero que indica la dirección de movimiento vertical. 1              = ascender el griper (potencia positiva en DC1). Cualquier otro = descender el griper (potencia negativa en DC1). \\
\end{tabular}\\[4pt]
\lstinputlisting[firstline=289,lastline=298,firstnumber=289]{../../../software/robot/main/Robot.py}
\subsection{\texttt{ChecarServo}}
\label{fn:ChecarServo}
Verifica el estado del smartservo 1 y lo detiene si detecta sobrecarga o bloqueo.

Espera 300 ms tras el inicio del movimiento para que el servo alcance velocidad estable, luego lee la corriente y la velocidad del smartservo\_1. Dependiendo del @p modo, aplica uno de dos criterios de parada: - Modo 1 (apertura): detiene el servo si la corriente supera 900 mA o la velocidad cae por debajo de 2 rpm, lo que indica que las pinzas llegaron al tope mecánico o encontraron resistencia excesiva. - Modo 0 (cierre): detiene el servo si la corriente supera 900 mA o la velocidad sube por encima de -2 rpm (el signo negativo refleja la dirección inversa del cierre), indicando bloqueo o tope mecánico. Al detectar la condición de parada, envía velocidad 0 al servo y espera 500 ms para que el mecanismo se estabilice antes de regresar el control.

\begin{tabular}{@{}p{0.15\linewidth}p{0.78\linewidth}@{}}
\textbf{Par\'ametro} & \textbf{Descripci\'on} \\
\texttt{modo} & Entero que indica la dirección de movimiento que se está supervisando. 1              = supervisión de apertura de pinzas. Cualquier otro = supervisión de cierre de pinzas. \\
\end{tabular}\\[4pt]
\lstinputlisting[firstline=317,lastline=330,firstnumber=317]{../../../software/robot/main/Robot.py}
\subsection{\texttt{Apertura\_Griper}}
\label{fn:AperturaGriper}
Controla la apertura o el cierre de las pinzas del griper.

Establece la potencia de aper\_pot\_griper en 50 \% y la aplica al canal DC2 de la placa de expansión de potencia. El signo positivo abre las pinzas y el negativo las cierra. La variable global aper\_pot\_griper se actualiza en cada llamada.

\begin{tabular}{@{}p{0.15\linewidth}p{0.78\linewidth}@{}}
\textbf{Par\'ametro} & \textbf{Descripci\'on} \\
\texttt{modo} & Entero que indica la acción a realizar sobre las pinzas. 1              = abrir las pinzas (potencia positiva en DC2). Cualquier otro = cerrar las pinzas (potencia negativa en DC2). \\
\end{tabular}\\[4pt]
Llama a: \texttt{ChecarServo} (\ref{fn:ChecarServo})\\
\lstinputlisting[firstline=342,lastline=352,firstnumber=342]{../../../software/robot/main/Robot.py}
\subsection{\texttt{Angulo\_Brush}}
\label{fn:AnguloBrush}
Ajusta el ángulo del smartservo 2 (brush) mediante el gamepad.

Incrementa o decrementa la variable global @p ang\_pot\_brush en pasos de 5 grados mientras la tecla correspondiente permanezca presionada, respetando los límites físicos del mecanismo: - Límite superior: 176 grados (modo 1, tecla "Left"). - Límite inferior: 4 grados   (modo 0, tecla "Right"). Tras cada paso aplica el nuevo ángulo al servo con una velocidad de movimiento de 20 unidades y espera 300 ms para que el servo alcance la posición antes de leer de nuevo el gamepad, evitando saltos bruscos. La variable global @p ang\_pot\_brush conserva el ángulo entre llamadas, de modo que el brush retiene su posición al soltar la tecla.

\begin{tabular}{@{}p{0.15\linewidth}p{0.78\linewidth}@{}}
\textbf{Par\'ametro} & \textbf{Descripci\'on} \\
\texttt{modo} & Entero que indica la dirección de ajuste del ángulo. 1              = incrementar el ángulo (mover hacia 176°, tecla "Left"). Cualquier otro = decrementar el ángulo (mover hacia 4°, tecla "Right"). \\
\end{tabular}\\[4pt]
\lstinputlisting[firstline=370,lastline=383,firstnumber=370]{../../../software/robot/main/Robot.py}
\subsection{\texttt{Lanzador}}
\label{fn:Lanzador}
Bucle interactivo de control de los motores brushless del lanzador.

Permite al operador ajustar la potencia de lanzamiento en tiempo real mediante el gamepad mientras el robot está en movimiento. El bucle se ejecuta hasta que el operador presiona N3: - R2 incrementa pot\_lanzamiento en 5 unidades (máximo 100). - L2 decrementa pot\_lanzamiento en 5 unidades (mínimo -100). - La potencia actualizada se aplica simultáneamente a BL1 y BL2. - N3 establece ambos motores a 0 y sale del bucle. La variable global pot\_lanzamiento conserva el último valor entre llamadas, por lo que la función lanzamiento() la pre-inicializa antes de llamar aquí.

\lstinputlisting[firstline=399,lastline=415,firstnumber=399]{../../../software/robot/main/Robot.py}
\subsection{\texttt{modo\_automatico}}
\label{fn:modoautomatico}
Activa el modo autónomo del robot durante un período de 28 segundos.

Reinicia el temporizador interno de novapi y ejecuta Comparacion\_QR() en un bucle continuo mientras el tiempo transcurrido sea menor a 28 s. Al expirar el tiempo, el bucle termina y el control regresa al llamador (normalmente el bucle principal de main()).

Llama a: \texttt{Comparacion\_QR} (\ref{fn:ComparacionQR})\\
\lstinputlisting[firstline=426,lastline=429,firstnumber=426]{../../../software/robot/main/Robot.py}
\subsection{\texttt{lanzamiento}}
\label{fn:lanzamiento}
Prepara e inicia el lanzador de pelotas con una potencia inicial del 20 \%.

Asigna 20 a la variable global pot\_lanzamiento como punto de partida seguro (suficiente para girar los motores sin un arranque brusco) y delega el control al bucle interactivo Lanzador().

Llama a: \texttt{Lanzador} (\ref{fn:Lanzador})\\
\lstinputlisting[firstline=437,lastline=440,firstnumber=437]{../../../software/robot/main/Robot.py}
\subsection{\texttt{Reset\_Motores}}
\label{fn:ResetMotores}
Detiene por completo los motores de tracción del robot.

Envía velocidad 0 a los cuatro motores encoder (M1–M4), deteniendo cualquier movimiento en curso. No afecta al griper ni al lanzador. Útil como parada de emergencia o al finalizar una secuencia autónoma.

\lstinputlisting[firstline=450,lastline=454,firstnumber=450]{../../../software/robot/main/Robot.py}
\subsection{\texttt{get\_pressed\_key}}
\label{fn:getpressedkey}
Detecta y devuelve la primera tecla registrada que esté presionada.

Itera sobre la tupla TECLAS en orden de declaración y consulta el estado de cada tecla al gamepad. Retorna en cuanto encuentra una tecla activa, garantizando un costo O(n) en el peor caso pero O(1) en promedio para la tecla más común. Si ninguna tecla está activa, retorna None para que el llamador no ejecute ninguna acción.

\lstinputlisting[firstline=465,lastline=469,firstnumber=465]{../../../software/robot/main/Robot.py}
\subsection{\texttt{main}}
\label{fn:main}
Ciclo principal del control del robot.

Ejecuta indefinidamente las siguientes operaciones en cada iteración: 1. Pausa de 1 ms (time.sleep(0.001)) para evitar rebotes de señal y reducir la carga de CPU al mínimo necesario. 2. Consulta get\_pressed\_key() para detectar teclas activas del gamepad. Si se detecta una tecla registrada en acciones\_gamepad, se ejecuta su función asociada con los argumentos definidos. 3. Lee los tres ejes del gamepad (Ly, Lx, Rx) para calcular las velocidades omnidireccionales mediante Checar\_Velocidades() y las aplica a los motores con Vel\_Motores(). El control manual de joystick es continuo: opera en paralelo con las acciones de tecla, salvo cuando una acción bloquea el hilo (p. ej. modo\_automatico() o Lanzador() tienen sus propios bucles internos).

Llama a: \texttt{get\_pressed\_key} (\ref{fn:getpressedkey}), \texttt{Checar\_Velocidades} (\ref{fn:ChecarVelocidades}), \texttt{Vel\_Motores} (\ref{fn:VelMotores})\\
\lstinputlisting[firstline=518,lastline=533,firstnumber=518]{../../../software/robot/main/Robot.py}
